\section{Reflection}\label{ch:reflection}

\subsection{From a broad question to a single test}\label{sec:refl-question}

The project began in Phase 1 with a broader question: does a regime-switching covariance model, in which volatilities and correlations both depend on a market state, forecast the tail risk of a stock--bond--gold portfolio better than a single-regime model? The pilot built the static models A to D on SPY, IEF and GLD and judged them against six go/no-go criteria that I wrote down before running them, covering correctness, the size and power of the correlation test, the stability of the estimated regimes, reproducibility and run time. All six passed, but two only narrowly: the power to detect a correlation breakdown reached 50\% only in the strongest simulated scenario, and the regimes were stable in 90.3\% of refits against a 90\% bar. I proceeded with the design, but these margins were the first sign that the correlation question would be hard to answer with daily data on three funds.

My supervisor then asked for a broad exploration of regime definitions and types of ETF. I split the sample so that 2005--2019 would serve for development and 2020--2025 would stay untouched until the test was fixed, and I logged every specification I tried. The exploration changed the question in two ways. First, the static comparison of regime-specific against pooled correlation, which the pilot had judged correctly sized, turned out to be oversized once volatility follows GARCH dynamics, so its many raw rejections were not evidence. This moved the question to models with GARCH margins, where the correlation channel could be isolated with the volatility model held fixed. Second, the long history of fund proxies showed that the development period contains no spell of positive stock--bond correlation, so only the hold-out could speak to the 2022 reversal. The question narrowed from ``do regimes help?'' to one comparison of one correlation, fixed in advance.

\subsection{Challenges and how I addressed them}\label{sec:refl-challenges}

The first challenge was theoretical framing. The literature disagrees about whether the stock--bond correlation follows a slow macroeconomic state or market stress, and my regime signal, built from equity volatility, can speak only to the second. I chose to state this openly as the reading the test addresses rather than to claim more, and Chapter~\ref{ch:discussion} reads the result against both.

The second was methodological. Comparing nested forecasting models out of sample is not covered cleanly by textbook tests, and the papers on the subject treat linear models under squared-error loss. Reading them led me to define the estimand as the expected loss difference of two forecasting procedures, in the sense of Giacomini and White, and to check the size of the test by simulation rather than rely on asymptotic theory. The power study then showed that six years of hold-out data would detect a moderate correlation shift only about one time in five. Rather than look for a more favourable test, I kept the one declared and added a smallest effect of interest and minimum detectable differences, so that a non-significant result would still say something.

The third concerned how results are worded. My supervisor's feedback on an earlier design document cautioned against overstating what each comparison isolates and against reading a non-significant result as evidence of no effect. I corrected that document, built Table~\ref{tab:differs} to state what each comparison changes, and fixed the wording of every possible outcome before the run.

\subsection{What the record shows and what I would change}\label{sec:refl-change}

Two episodes tested the discipline of the design. While searching the test files before the freeze, I displayed two single-day hold-out returns by accident; nothing depended on them, and I disclosed them. And after the hold-out forecasts were produced, the scorer refused them because they had been recorded two commits after the frozen tag, although the forecast code was identical. I amended the check before any number was computed and recorded the change in the specification. Both episodes showed me that a pre-registration kept inside one's own repository is only as credible as the record around it, which is why the appendix sets them out in full.

With hindsight I would change three things. I would include the practitioner baselines, historical simulation and an EWMA, in the frozen design, since the simplest of them turned out to be the best calibrated. I would choose the regime signal from the economic argument, for example from inflation or from the stock--bond correlation itself, rather than from the exploration and the power study. And I would treat calibration as a primary outcome alongside the loss comparison, because for a risk desk it was the more consequential result. What I would keep is the order of work: explore freely, write the test down, then use the hold-out once.
